\chapter{Correctness and Invariants}

\section{Testing Philosophy}

The test suite is designed to check the public API across all backends.
Each backend is run through the same scenarios, so failures reveal either backend-specific bugs or inconsistencies in the common API contract.

The tests use a custom \texttt{TEST\_ASSERT} macro rather than the standard \texttt{assert}.
This keeps tests active even in optimized \texttt{-DNDEBUG} builds.

\section{Deterministic Tests}

The deterministic tests cover:

\begin{itemize}
    \item rejecting missing comparators;
    \item behavior of empty heaps;
    \item insertion and sorted extraction;
    \item duplicate priorities;
    \item interleaved insertion and extraction;
    \item larger deterministic ordering;
    \item \texttt{contains} and remove;
    \item remove by handle identity rather than by priority equality;
    \item stale handle rejection;
    \item wrong-heap handle rejection;
    \item decrease-key behavior;
    \item targeted operations after backend restructuring.
\end{itemize}

These cases are important because several bugs in heap implementations are not visible through simple insert/extract tests.
Handles, restructuring, and stale references are more demanding.

\section{Randomized Oracle Test}

The randomized test maintains a simple external oracle.
It randomly performs inserts, decrease-key operations, removes, and extract-min operations.
After each step, it compares heap state with the oracle:

\begin{itemize}
    \item heap size matches active item count;
    \item empty state matches active item count;
    \item peek-min matches the minimum active item;
    \item extracted items match oracle minima.
\end{itemize}

The default seed is deterministic.
The seed, number of items, and number of steps can be overridden:

\begin{lstlisting}[style=shellstyle, caption={Longer randomized campaign.}]
HEAPX_TEST_SEED=123 \
HEAPX_TEST_ITEMS=512 \
HEAPX_TEST_STEPS=20000 \
make test
\end{lstlisting}

\section{Common Invariants}

\texttt{heapx\_check\_invariants()} first validates common state:

\begin{itemize}
    \item heap pointer is non-null;
    \item vtable and comparator are present;
    \item heap id is non-zero;
    \item handle slot count does not exceed capacity;
    \item handle table pointer and capacity agree;
    \item free-list links are valid and acyclic within the slot count;
    \item free slots do not claim live owners or items.
\end{itemize}

After the common checks, the function dispatches to the backend invariant checker.

\section{Binary Heap Invariants}

The binary heap checker verifies:

\begin{itemize}
    \item size does not exceed capacity;
    \item array pointer and capacity agree;
    \item every locator points back to its current array index;
    \item every locator resolves through the public handle table;
    \item every parent compares no worse than its children.
\end{itemize}

This checks both algorithmic heap order and handle consistency.

\section{Fibonacci Heap Invariants}

The Fibonacci checker recursively validates circular root and child lists.
It checks:

\begin{itemize}
    \item empty heaps have no minimum pointer;
    \item non-empty heaps have a minimum pointer;
    \item roots have no parent;
    \item the minimum pointer is not worse than any root;
    \item sibling links are bidirectionally consistent;
    \item child parent pointers are correct;
    \item child count matches node degree;
    \item child nodes do not violate heap order with their parent;
    \item handled nodes resolve to themselves;
    \item counted nodes match heap size.
\end{itemize}

\section{Kaplan Heap Invariants}

The Kaplan checker validates a different structure:

\begin{itemize}
    \item reusable rank table entries are clear between consolidations;
    \item empty heaps have no root;
    \item non-empty heaps have exactly one root tree;
    \item root parent and sibling links are null;
    \item child and sibling links are consistent;
    \item mark bits are valid;
    \item ranks do not exceed heap size or direct child count;
    \item heap order holds between parent and child;
    \item handled nodes resolve to themselves;
    \item counted nodes match heap size.
\end{itemize}

These checks are especially valuable because the Kaplan backend uses rank-state repairs and temporary root lists during extract-min.

\section{Sanitizers and CI}

The Makefile provides several verification targets:

\begin{lstlisting}[style=shellstyle, caption={Verification targets.}]
make test
make debug-checks
make sanitize
make leak-check
make optimized-check
make ci
\end{lstlisting}

\begin{itemize}
    \item \texttt{make sanitize} enables AddressSanitizer and UndefinedBehaviorSanitizer.
    \item \texttt{make leak-check} runs leak detection where supported.
    \item \texttt{make optimized-check} tests with \texttt{-O2 -DNDEBUG}.
    \item The CI workflow runs \texttt{make ci} and checks whitespace with \texttt{git diff --check}.
\end{itemize}
